% ************************** Thesis Abstract *****************************
% Use `abstract' as an option in the document class to print only the titlepage and the abstract.
\begin{abstract}

To meet the demands of increased luminosity at the High-Luminosity Large Hadron Collider (HL-LHC) in 2036, new generations of silicon detectors must achieve higher spatial and timing resolutions and improved radiation tolerance. As part of the LHCb Upgrade II, the Mighty Tracker system will implement High-Voltage Monolithic Active Pixel Sensors (HV-MAPS) based on Complementary Metal-Oxide-Semiconductor (CMOS) technology. MightyPix, a candidate HV-MAPS sensor for this upgrade, is currently in limited supply. Therefore, this project investigates the performance of TelePix2 - a sister sensor with similar architecture and broader availability.

The experiment was performed at the DESY II facility in March, 2024. Test beam data were collected with a setup that included three telescope planes, a scintillator trigger, and a 4 GeV electron beam. Event reconstruction was performed using the Corryvreckan framework.

Time-over-threshold (ToT) and cluster size distributions were analysed across runs with threshold from 140-180 DAC and bias voltages 40-85 V on pixel, 120-130 V on chip. The ToT peak increased with bias voltage and decreased with threshold, consistent with expectations. The mean cluster size exhibits a minimum that shifts to higher threshold values for increasing bias voltage. This behaviour can be explained by the interplay between sensor characteristics and the readout system.

Based on efficiency analysis, the optimal configuration was identified, and under these conditions (130 V chip bias, 140 DAC threshold), the time resolution of TelePix2 was measured to be $2.656 \pm 0.003$ ns.

\end{abstract}






